\section*{Introduction}

Modal logic is most commonly interpreted over Kripke frames, 
but it also admits a variety of spatial semantics, 
in which formulas are read as denoting 
regions of a topological or geometric space.
One of the earliest results of this kind is the McKinsey--Tarski 
interpretation of $S4$ over arbitrary topological spaces, 
with the diamond modality read as topological closure. 
A recurring theme in this line of work is to ask how much 
of the classical model theory of modal logic 
(soundness and completeness, but also correspondence-theoretic results)
survives the move from frames to spaces.

One of the cornerstones of correspondence theory for Kripke frames
is the Goldblatt-Thomason theorem: 
a class of frames is modally definable 
if and only if it is closed under finite disjoint unions, 
generated subframes, p-morphic images, and it's closed under 
and reflects ultrafilter extensions
\cite{Blackburn_Rijke_Venema_2001}. 
This result gives a purely semantic, 
structural characterization of the classes of frames 
that can be picked out by a set of modal formulas, 
and analogues of it have since been sought 
in other semantics for modal logic. 
In the topological setting, ten Cate, Gabelaia and Sustretov 
obtained such an analogue for modal languages 
interpreted over topological spaces \cite{TENCATE}, 
building on earlier work on modal definability in topology 
\cite{Gabelaia2001ModalDI}.

More recently, attention has turned to a semantics 
that sits between the relational and the topological: 
polyhedral semantics, in which formulas are evaluated on 
compact polyhedra, with the diamond of a formula 
interpreted as the closure of the subpolyhedron to which the formula evaluates. 
This semantics was developed systematically by Dekker \cite{Dekker}, 
and is closely tied to applications in spatial 
and geometric model checking \cite{DBLP}. 

This new semantics raises a natural question, 
which is the starting point of the present paper: 
does a Goldblatt-Thomason-type theorem hold for compact polyhedra? 
Answering it is not simply a matter of transporting the classical 
statement across the correspondence. 
Compact polyhedra do not admit generated subpolyhedra 
in the naive sense: the polyhedral realization of a 
generated subframe is an open subpolyhedron, which in general fails to be 
compact, and formulating the equivalent notion of a p-morphism in this setting
is a non-trivial task.

Our main contribution is such a reformulation, 
together with a proof that it is exactly the right one. 
We introduce the notion of a \textbf{face up-reduction} 
between compact polyhedra: 
a partial map into the face poset of a triangulation of the target polyhedron, 
required to be open and surjective, 
and to pull back opens of the Alexandrov topology to open subpolyhedra. 
Face up-reductions play the role that p-morphic images 
of generated subframes play classically. Our central result, Theorem 3.2.1 (GTKP), 
states that a class of compact polyhedra is modally definable 
if and only if it is closed under finite disjoint unions and 
face up-reductions.

To prove this we develop, largely following \cite{Dekker}, 
the combinatorics of triangulations needed to move between 
a polyhedron and any one of its face posets: 
barycentric subdivision, common triangulations of 
a polyhedron and a subpolyhedron, 
and the fact that repeated subdivision of a triangulation 
eventually refines any other triangulation of the same polyhedron. 
These facts, standard in piecewise-linear topology \cite{PLTopology}, 
let us transport p-morphisms and Jankov-Fine formulas 
between a polyhedron and its triangulations without changing 
the logic being described.

The paper is organized as follows:
\begin{itemize}
    \item Section 1 recalls simplicial complexes, polyhedra, 
and the correspondence between simplicial complexes, 
abstract simplicial complexes, and posets via the nerve construction, 
together with the subdivision lemmas used throughout the paper. 
    \item Section 2 sets up polyhedral semantics, 
its connection to the relational semantics of $S4.\text{Grz}$,
and recalls the classical Goldblatt-Thomason theorem and 
the Jankov-Fine axiomatization in the form we will need.
    \item Section 3 introduces face up-reductions and proves our main theorem, 
along with a corollary for simplicial up-reductions and 
a lemma showing the result does not depend on the choice of triangulation. 
    \item Section 4 provides concrete examples of modally definable 
and non-modally-definable classes of compact polyhedra.
    \item Section 5 closes the paper discussing open questions, 
chiefly a conjectured strengthening of our theorem 
to simplicial up-reductions, 
and examines what changes if we allow open or infinite polyhedra. 
\end{itemize}

 




